\documentclass[12pt]{article}

\usepackage{projeto}
\usepackage{graphicx,url}
\usepackage[utf8]{inputenc}
\usepackage[brazil]{babel}

     
\sloppy

\title{CMP263 - Machine Learning Paper \\ Japanese House Prices Predition}

\author{Juliano Almeida Machado\inst{1}}

\address{Instituto de Informática -- Universidade Federal do Rio Grande do Sul (UFRGS)\\
  Caixa Postal 15.064 -- 91.501-970 -- Porto Alegre -- RS -- Brazil}

\begin{document} 

\maketitle

\begin{abstract}
  This meta-paper describes the style to be used in articles and short papers for SBC conferences. For papers in
  English, you should add just an abstract while for the papers in Portuguese, we also ask for an abstract in Portuguese
  (``resumo''). In both cases, abstracts should not have more than 10 lines and must be in the first page of the paper.
\end{abstract}
     
\begin{resumo} 
  Este meta-artigo descreve o estilo a ser usado na confecção de artigos e resumos de artigos para publicação nos anais
  das conferências organizadas pela SBC. É solicitada a escrita de resumo e abstract apenas para os artigos escritos em
  português. Artigos em inglês deverão apresentar apenas abstract. Nos dois casos, o autor deve tomar cuidado para que o
  resumo (e o abstract) não ultrapassem 10 linhas cada, sendo que ambos devem estar na primeira página do artigo.
\end{resumo}


\section{Introduction}
The Japanese housing market presents a anomaly within global real estate economics. In most developed nations,
residential property functions as an appreciating asset and a means for generational wealth accumulation. In contrast,
housing in Japan behaves as a depreciating consumer durable. Within 20 to 30 years of construction, the structural value
of a typical Japanese home depreciates to virtually zero. This structural depreciation creates a stark institutional
decoupling between land value and building value. While urban land, especially within the Tokyo and Osaka metropolitan
areas, retains substantial value and responds to conventional market forces, the structures erected upon that land are
treated as temporary assets. According to data from the Ministry of Land, Infrastructure, Transport and Tourism (MLIT),
the secondary housing market accounts for only approximately 15\% of total housing transactions in Japan, compared to
over 80\% to 90\% in the US and UK. 

This environment traces its origins back to the post-World War II reconstruction era. Facing severe housing shortages
and rapid urbanization, the Japanese government prioritized speed and quantity over long-term structural durability. The
resulting regulatory framework and market incentives codified a system of high-turnover construction, which was later
exacerbated by the asset price bubble of the late 1980s. When the bubble burst in 1991, it left a legacy of cautious
institutional lending and a deeply fragmented real estate framework.

% #TODO: Check if this is true in the data; 
Japan’s geographic vulnerability to seismic activity also dictates its building philosophies. The Building Standard Act
of Japan has undergone continuous revisions, with the most critical shift occurring in 1981 with the introduction of the
\textit{Shin-Taishin} (New Anti-Seismic Design Code). Homes constructed prior to 1981 are fundamentally viewed as
structurally unsafe and financially uninsurable by contemporary standards. Subsequent amendments following the 1995
Hanshin-Awaji (Kobe) earthquake and the 2011 Tohoku earthquake have further accelerated the technological obsolescence
of older properties. Consequently, buyers often perceive purchasing a pre-owned home not just as an economic compromise,
but as a direct compromise on structural safety.

Complementing these rigid safety mandates is a pervasive cultural preference for new builds, known as
\textit{Shinchiku}. A newly built home carries a significant socio-cultural premium in Japan, symbolizing clean slates,
modernization, and social status. This preference is structurally reinforced by financial institutions, whereas Japanese
banks routinely structure residential mortgages around the statutory useful life (\textit{Taikyusei}) of the building.
For a standard timber-frame home, this period is legally defined as 22 years. Once a structure surpasses this threshold,
securing a mortgage for its resale becomes exceedingly difficult, effectively freezing the secondary market.

It is also important to note that one of the most pressing modern challenge facing the Japanese housing market is the
nation's demographic trajectory. Japan is characterized by a super-aging society combined with a sub-replacement
fertility rate. The population has been in a continuous state of decline since its peak in the late 2000s. This
population contraction has triggered the abandoned house (\textit{akiya}) phenomenon. Millions of homes across rural,
suburban, and even peripheral urban areas stand completely vacant, given the tendency of inheriting families to
frequently abandon the physical structures. This has created a massive inventory of valueless, deteriorating properties
that distort local real estate valuations and create severe municipal management challenges.

The Japanese housing market presents an invaluable case study for urban economists and sociologists alike. It challenges
the universal assumption that residential property is an inherently safe, appreciating asset class. As other
post-industrial nations begin to face similar demographic plateaus and climate-induced regulatory pressures,
understanding the mechanisms of Japan’s ``scrap-and-build'' housing model becomes essential. 

\section{Problem}

\bibliographystyle{projeto}
\bibliography{projeto}

\end{document}